%Section liste des commandes

\part{Liste des commandes, par thème}

\begin{noteblock}
\cmaj{2.0.0} Cette section contient un \textit{résumé} des différentes commandes et environnements disponibles dans \ctex{ProfLycee}.

Elles sont présentées de manière \textit{succincte}, mais elles sont présentées de manière \textit{détaillée} dans la suite de la documentation.

Autant que peut se faire, la version \textit{alias} en \ctex{\textbackslash pfl...} est mentionnée, en sachant qu'elle reprend la syntaxe complète de la commande initiale.
\end{noteblock}

\begin{PresCodeTexPL}{listing only}
%calculs formatés
\pflnum(*)<type>[dec]{calcul}                      %version consolidée
\pflnumsqrt[formatage]{intérieur de la racine}
\pflnumfrac(*)[formatage]{calcul}
\pflnumrad(*)[formatage]{calcul}
\end{PresCodeTexPL}

\begin{PresCodeTexPL}{listing only}
%conversion directe entre bases
\pflbasetobase{base initiale}{base finale}{nb}
\end{PresCodeTexPL}

\begin{PresCodeTexPL}{listing only}
%formatage spéciaux avec sinuitx
\pflprix[arrondi]{px}
\pflangle[arrondi]{agl}
\pflpcent[arrondi]{pcent}
\end{PresCodeTexPL}

\begin{PresCodeTexPL}{listing only}
%intervalles
\begin{RepIntervalles}[clés]<options tikz>...\end{RepIntervalles}
\tkzIntervalle[clés]{xmin}{xmax}
\end{PresCodeTexPL}

\begin{PresCodeTexPL}{listing only}
%Résolution approchée d'une équation f(x)=k
\ResolutionApprochee[clés]{équation}[macro]
\pflresolapproch
\end{PresCodeTexPL}

\begin{PresCodeTexPL}{listing only}
%Présentation d'une solution par balayage (TVI)
\SolutionTVI[options]{fonction}{valeur}
\pflsoltvi
\end{PresCodeTexPL}

\begin{PresCodeTexPL}{listing only}
%Calculer le terme d'une suite récurrente simple, toile pour une suite récurrente simple
\CalculTermeRecurrence[options]{fonction associée}
\pflcalcrecurr
\ToileRecurrence[clés][options du tracé][option supplémentaire des termes]
\pfltoilerecurr
\end{PresCodeTexPL}

\begin{PresCodeTexPL}{listing only}
%Mise en forme de la conclusion d'un seuil
\SolutionSeuil[options]{fonction associée}{seuil}
\pflsolseuil
\end{PresCodeTexPL}

\begin{PresCodeTexPL}{listing only}
%Valeur approchée d'une intégrale
\IntegraleApprochee[clés]{fonction}{a}{b}
\pflintegrapproch
\CalcIntegrale(*)[prec]<nb subdiv>{expr version xint}{binf}{bsup}
\pflvalintegr
\end{PresCodeTexPL}

\begin{PresCodeTexPL}{listing only}
%Valeur moyenne d'une fonction
\ValeurMoyenneIntg(*)[prec]<nb subdiv>{expr version xint}{binf}{bsup}
\pflvalmoyintegr
\end{PresCodeTexPL}

\begin{PresCodeTexPL}{listing only}
%forme canonique
\FormeCanonique(*)[option]{a}{b}{c}
\pflformcanoniq
\end{PresCodeTexPL}

\begin{PresCodeTexPL}{listing only}
%décomposition d'une fonction homographique
\FonctionHomographique(*){a}{b}{c}{d}
\pflfcthomogr
\end{PresCodeTexPL}

\begin{PresCodeTexPL}{listing only}
%max/min d'une fonction
\DetermineMax[précision]{expression xint}{a}{b}[\tmpmax][\tmpmaxvalx]
\pfldetmax
\DetermineMin[précision]{expression xint}{a}{b}[\tmpmin][\tmpminvalx]
\pfldetmin
\end{PresCodeTexPL}

%\begin{PresCodeTexPL}{listing only}
%fenêtre de repérage en tikz et courbe
%\GrilleTikz[options][options grille ppale][options grille second.]
%\AxesTikz[options] \AxexyTikz[optionsOx][optionsOy]{valeursOx}{valeursOy}
%\AxexTikz[options]{valeurs} \AxeyTikz[options]{valeurs}
%
%\FenetreSimpleTikz[options](opt axes)<opt axe Ox>{liste valx}<opt axe Oy>{liste valy}
%\DeclareFonctionTikz[nom]{expr}
%\DeclareFonctionTikzXint[nom]{expr}
%\CourbeTikz[options]{fonction}{valxmin:valxmax}
%\CourbeTikzXint[options tikz]{fonction xint}{deb..[pas]..fin}
%
%%génération du tracé d'une courbe d'interpolation, dans une commande tikz
%\GenereSplineTikz[options]{liste}[\nomdutracé]
%
%%courbe d'interpolation, tangente, dans un environnement tikz
%\SplineTikz[options]{liste}
%\TangenteTikz[options]{liste}
%\PtsDiscontinuite[options]{liste}

\begin{PresCodeTexPL}{listing only}
%schémas pour le signe affine/trinôme, dans un environnement tikz
\MiniSchemaSignes(*)[clés]<options tikz>
\pflschemasignes
\MiniSchemaSignesTkzTab[options]{numligne}[échelle][décalage horizontal]
\pflschemasignestkztab
\end{PresCodeTexPL}

\begin{PresCodeTexPL}{listing only}
%intégrales et méthodes graphiques
\IntegraleApprocheeTikz[clés]{nom_fonction}{a}{b}
\pflintegrapprochtikz
\end{PresCodeTexPL}

\begin{PresCodeTexPL}{listing only}
%présentation de code Python
\begin{CodePythonLst}(*)[clés]{commandes tcbox}...\end{CodePythonLst}
\begin{CodePythonLstAlt}(*)[clés]{commandes tcbox}...\end{CodePythonLstAlt}
%:=librairie piton
\begin{CodePiton}[options piton]{commandes tcbox}<option line-numbers>...\end{CodePiton}
\begin{PitonConsole}<clés>{commandes tcbox}...\end{PitonConsole}
%:=librairie pythontex
\begin{CodePythontex}[clés]{commandes tcbox}...\end{CodePythontex}
\begin{CodePythontexAlt}[clés]{commandes tcbox}...\end{CodePythontexAlt}
\begin{ConsolePythontex}[options]{}...\end{ConsolePythontex}
%présentation de pseudocode
\begin{PseudoCode}(*)[clés]{commandes tcbox}...\end{PseudoCode}
\begin{PseudoCodeAlt}(*)[largeur][clés]{commandes tcbox}...\end{PseudoCodeAlt}
\begin{PseudoCodePiton}[clés]{commandes tcbox}<option line-numbers>...\end{PseudoCodePiton}

%style Thonny
\begin{PitonThonnyEditor}<clés>[options tcbox]{largeur}...\end{PitonThonnyEditor}
\begin{PitonThonnyConsole}<clés>[options tcbox]{largeur}...\end{PitonThonnyConsole}

\CartoucheCapytale(*)[options]{code capytale}
\end{PresCodeTexPL}

\begin{PresCodeTexPL}{listing only}
%présentation de codes via piton (insertion de fichiers)
\CodePitonFichier[clés]{fichier}<option line-numbers>
\PseudoCodePitonFichier[clés]{fichier}<option line-numbers>
\PitonThonnyEditorFichier<clés>[largeur]{fichier}
\end{PresCodeTexPL}

\begin{PresCodeTexPL}{listing only}
%pavé et tétraèdre, dans un environnement tikz
\PaveTikz[options]
\TetraedreTikz[options]
\end{PresCodeTexPL}

\begin{PresCodeTexPL}{listing only}
%cercle trigo, dans un environnement tikz
\CercleTrigo[clés]
\pflcercletrigo
\end{PresCodeTexPL}

\begin{PresCodeTexPL}{listing only}
%Affichage des coordonnées d'un point (2 ou 3 coordonnées)
\AffPoint[options de formatage](liste des coordonnées)
%Affichage des coordonnées d'un vecteur (2 ou 3 coordonnées)
\AffVecteur[options de formatage]<options nicematrix>(liste des coordonnées)
\end{PresCodeTexPL}

\begin{PresCodeTexPL}{listing only}
%Avec un vecteur normal et un point
\TrouveEqCartPlan[clés](vecteur normal)(point)
%Avec deux vecteurs directeurs et un point
\TrouveEqCartPlan[clés](vecteur dir1)(vecteur dir2)(point)
%Avec trois points
\TrouveEqCartPlan[clés](point1)(point2)(point3)
\end{PresCodeTexPL}

\begin{PresCodeTexPL}{listing only}
%Avec un vecteur directeur et un point
\TrouveEqParamDroite[clés](vecteur directeur)(point)
%Avec deux points
\TrouveEqParamDroite[clés](point1)(point2)
\end{PresCodeTexPL}

\begin{PresCodeTexPL}{listing only}
%Avec un vecteur normal (choix par défaut) et un point
\TrouveEqCartDroite[clés](vecteur normal)(point)
%Avec un vecteur directeur et un point
\TrouveEqCartDroite[clés,VectDirecteur](vecteur directeur)(point1)
%Avec deux points
\TrouveEqCartDroite[clés](point1)(point2)
\end{PresCodeTexPL}

\begin{PresCodeTexPL}{listing only}
%Avec le point et le plan via vect normal + point
\TrouveDistancePtPlan(point)(vec normal du plan)(point du plan)
%Avec le point et le plan via vect normal + point
\TrouveDistancePtPlan(point)(équation cartésienne)
\end{PresCodeTexPL}

\begin{PresCodeTexPL}{listing only}
\TrouveEqSphere[clés](centre ou point n°1)(rayon ou point ou éq plan)
\end{PresCodeTexPL}

\begin{PresCodeTexPL}{listing only}
%Avec le vecteur
\TrouveNorme(vecteur)
%Avec deux points
\TrouveNorme(point 1)(point 2)
\end{PresCodeTexPL}

\begin{PresCodeTexPL}{listing only}
%Équation réduite d'une droite
\EquationReduite[option]{A/xa/ya,B/xb/yb}
\pfleqreduite
\end{PresCodeTexPL}

\begin{PresCodeTexPL}{listing only}
%Schémas de cours de géométrie dans l'espace
\SchemaEspace(*)[clés]{type}
\pflschemesp
\end{PresCodeTexPL}

\begin{PresCodeTexPL}{listing only}
%paramètres d'une régression linéaire, nuage de points
\CalculsRegLin[clés]{listeX}{listeY}
\PointsRegLin[clés]{listeX}{listeY}
\pflcalcreglin
\end{PresCodeTexPL}

\begin{PresCodeTexPL}{listing only}
%stats à 2 variables, dans un environnement tikz
\GrilleTikz[options][options grille ppale][options grille second.]
\AxesTikz[options]
\AxexTikz[options]{valeurs} \AxeyTikz[options]{valeurs}
\FenetreTikz \OrigineTikz
\FenetreSimpleTikz[options](opt axes)<opt axe Ox>{liste valx}<opt axe Oy>{liste valy}
\NuagePointsTikz[options]{listeX}{listeY}
\PointMoyenTikz[options]
\CourbeTikz[options]{formule}{domaine}
\end{PresCodeTexPL}

\begin{PresCodeTexPL}{listing only}
%paramètres d'une série discrète
\DeterminerParamStats(*){liste}[\macromin]...[\macromax]
\pflparamstats
\end{PresCodeTexPL}

\begin{PresCodeTexPL}{listing only}
%boîte à moustaches, dans un environnement tikz
\BoiteMoustaches[options]
\BoiteMoustachesAxe[options]
\pflboitemoustach
\end{PresCodeTexPL}

\begin{PresCodeTexPL}{listing only}
%histogrammes
\Histogramme(*)[options]{données}
\pflhistogramme
\end{PresCodeTexPL}

\begin{PresCodeTexPL}{listing only}
%courbe ECC/FCC
\CourbeECC[clés]{liste valeurs}{liste effectifs}
\pflcourbeecc

\begin{EnvCourbeECC}[clés]{liste valeurs}{liste effectifs}...\end{EnvCourbeECC}
\MedianeQuartilesECC{liste valeurs}{liste effectifs}
\end{PresCodeTexPL}

\begin{PresCodeTexPL}{listing only}
%loi binomiale B(n,p)
\CalcBinomP{n}{p}{k}
\CalcBinomC{n}{p}{a}{b}
\BinomP(*)[prec]{n}{p}{k}
\BinomC(*)[prec]{n}{p}{a}{b}
\pflbinomp
\pflbinomc
\end{PresCodeTexPL}

\begin{PresCodeTexPL}{listing only}
%loi de Poisson P(l)
\CalcPoissP{l}{k}
\CalcPoissC{l}{a}{b}
\PoissonP(*)[prec]{l}{k}
\PoissonC(*)[prec]{l}{a}{b}
\pflpoissonp
\pflpoissonc
\end{PresCodeTexPL}

\begin{PresCodeTexPL}{listing only}
%loi géométrique G(p)
\CalcGeomP{p}{k}
\CalcGeomC{l}{a}{b}
\GeomP{p}{k}
\GeomC{l}{a}{b}
\pflgeomp
\pflgeomc
\end{PresCodeTexPL}

\begin{PresCodeTexPL}{listing only}
%loi hypergéométrique H(N,n,m)
\CalcHypergeomP{N}{n}{m}{k}
\CalcHypergeomP{N}{n}{m}{a}{b}
\HypergeomP{N}{n}{m}{k}
\HypergeomC{N}{n}{m}{a}{b}
\pflhypergeomp
\pflhypergeomc
\end{PresCodeTexPL}

\begin{PresCodeTexPL}{listing only}
%loi normale N(m,s)
\CalcNormC{m}{s}{a}{b}
\NormaleC(*)[prec]{m}{s}{a}{b}
\pflnormalec
\end{PresCodeTexPL}

\begin{PresCodeTexPL}{listing only}
%loi exponentielle E(l)
\CalcExpoC{l}{a}{b}
\ExpoC(*)[prec]{l}{a}{b}
\pflexpoc
\end{PresCodeTexPL}

\begin{PresCodeTexPL}{listing only}
%arbres de probas
\ArbreProbasTikz[options]<opts tikz>{donnees}
\pflarbreprobas

\begin{EnvArbreProbasTikz}[options]<opts tikz>{donnees}...\end{EnvArbreProbasTikz}
\end{PresCodeTexPL}

\begin{PresCodeTexPL}{listing only}
%schémas lois continues
\LoiNormaleGraphe[options]<options tikz>{m}{s}{a}{b}
\LoiExpoGraphe[options]<options tikz>{l}{a}{b}
\pflnormgraph
\pflexpograph
\end{PresCodeTexPL}

\begin{PresCodeTexPL}{listing only}
%fonction de répartition discrète, dans une environnement tikz
\FonctionRepartTikz[clés]{probas,borneinf,bornesup / probas,borneinf,bornesup / ...}
\end{PresCodeTexPL}

\begin{PresCodeTexPL}{listing only}
%binomiale
\HistogrammeBinomiale[clés]{n}{p}
\pflhistobinom

\begin{HistoBinomiale}[clés]{n}{p}...\end{HistoBinomiale}

\TraceLoiNormale(*)[clés]<plage>{param1}{param2}
\end{PresCodeTexPL}

\begin{PresCodeTexPL}{listing only}
%entier aléatoire entre a et b
\NbAlea{a}{b}{macro}
\pflnbalea
\end{PresCodeTexPL}

\begin{PresCodeTexPL}{listing only}
%nombre décimal (n chiffres après la virgule) aléatoire entre a et b+1 (exclus)
\NbAlea[n]{a}{b}{macro}
\end{PresCodeTexPL}

\begin{PresCodeTexPL}{listing only}
%création d'un nombre aléatoire sous forme d'une macro
\VarNbAlea{macro}{calcul}
\pflvarnbalea
\end{PresCodeTexPL}

\begin{PresCodeTexPL}{listing only}
%liste d'entiers aléatoires
\TirageAleatoireEntiers[options]{macro}
\pfltiragealeaent
\end{PresCodeTexPL}

\begin{PresCodeTexPL}{listing only}
%liste des diviseurs
\ListeDiviseurs(*)[option]{nombre}
\pfllistediv
\end{PresCodeTexPL}

\begin{PresCodeTexPL}{listing only}
%arbre des diviseurs
\ArbreDiviseurs[options]{nombre}
\pflarbrediv
\end{PresCodeTexPL}

\begin{PresCodeTexPL}{listing only}
%test booléen nb premier
\pflboolestpremier{nb}
\pflestpremier[code si vrai][code si faux]{nb}
\end{PresCodeTexPL}

\begin{PresCodeTexPL}{listing only}
%décomposition avec exposants en nombres premiers
\pfldecompnb(*){nb}
\end{PresCodeTexPL}

\begin{PresCodeTexPL}{listing only}
%arbre "avec remise"
\ArbreChoix[clés]<options tikzpicture>{liste}
\pflarbrechoix
\end{PresCodeTexPL}

\begin{PresCodeTexPL}{listing only}
%arbre "sans remise"
\ArbreChoixSansRemise[clés]<options tikzpicture>{liste}
\pflarbrechoixssremise
\end{PresCodeTexPL}

\begin{PresCodeTexPL}{listing only}
%factorielle
\Factorielle(*)[clés]{n}
\pflfacto

%primorielle
\Primorielle(*)[clés]{n}
\pflprimor
\end{PresCodeTexPL}

\begin{PresCodeTexPL}{listing only}
%arrangement Anp
\Arrangement(*)[option]{p}{n}
\pflarrang
\end{PresCodeTexPL}

\begin{PresCodeTexPL}{listing only}
%arrangement Cnp (p parmi n)
\Combinaison(*)[option]{p}{n}
\pflcombin
\end{PresCodeTexPL}

\begin{PresCodeTexPL}{listing only}
%conversions
\ConversionEntreBases(*)[espace blocs 4 pour le binaire]{dep->arrivée}{nb}
\pflconvbases

\ConversionDecBin(*)[clés]{nombre}
\pflconvdecbin

\ConversionBinHex[clés]{nombre}
\pflconvbinhex

\ConversionHexBin[clés]{nombre}
\pflconvhexbin

\ConversionVersDec[clés]{nombre}
\pflconvversdec

\ConversionBaseDix[clés]{nombre}{base de départ}
\pflconvdix

\ConversionDepuisBaseDix[options]{nombre en base 10}{base d'arrivée}
\pflconvdepuisdix
\end{PresCodeTexPL}

\begin{PresCodeTexPL}{listing only}
%PGCD présenté
\PresentationPGCD[options]{a}{b}
\pflprespgcd
\end{PresCodeTexPL}

\begin{PresCodeTexPL}{listing only}
%Opérations posées
\OperationPosee[clés]{opération}
\pflopeposee
\end{PresCodeTexPL}

\begin{PresCodeTexPL}{listing only}
%Équation diophantienne
\EquationDiophantienne[clés]{equation}
\pflequadioph
\end{PresCodeTexPL}

\begin{PresCodeTexPL}{listing only}
%division euclidienne présentée
\DivEucl(*)[clés]{a}{b}
\pfldiveucl

%division euclidienne 'brute'
\DivisionEucl{a}{b}
\pflpresdiveucl
\end{PresCodeTexPL}

\begin{PresCodeTexPL}{listing only}
%reste et quotient modulo de a par n
\QuoMod(*){a}{n}
\pflquomod

\ResteMod(*){a}{n}
\pflrestemod
\end{PresCodeTexPL}

\begin{PresCodeTexPL}{listing only}
%inverse modulo
\InverseModulo(*){a}{modulo}
\pflinvmod
\end{PresCodeTexPL}

\begin{PresCodeTexPL}{listing only}
%Chiffrement de César
\ChiffrementCesar[clés]{message}
\pflchiffrcesar
\end{PresCodeTexPL}

\begin{PresCodeTexPL}{listing only}
%Chiffrement affine
\ChiffrementAffine[clés]{message}
\pflchiffraff
\end{PresCodeTexPL}

\begin{PresCodeTexPL}{listing only}
%Chiffrement de Hill
\ChiffrementHill[clés]{message}
\pflchiffrhill
\end{PresCodeTexPL}

\begin{PresCodeTexPL}{listing only}
%conversion en fraction, simplification de racine
\ConversionFraction(*)[option]{argument}
\SimplificationRacine{expression}
\end{PresCodeTexPL}

\begin{PresCodeTexPL}{listing only}
%ensemble d'éléments
\EcritureEnsemble[clés]{liste}
\pflecritens
\end{PresCodeTexPL}

\begin{PresCodeTexPL}{listing only}
%trinôme, trinôme aléatoire
\EcritureTrinome[options]{a}{b}{c}
\pflecrittrinom
\end{PresCodeTexPL}

\begin{PresCodeTexPL}{listing only}
%mesure principale
\MesurePrincipale[options]{angle}
\pflmesppale

%lignes trigo
\LigneTrigo(*)[booléens]{cos/sin/tan}(angle)
\pfllignetrig
\end{PresCodeTexPL}

\begin{PresCodeTexPL}{listing only}
%compétences LEGT
\CompMathsLyc[clés]{choix}
%compétences LPRO
\CompMathsLycPro[clés]{choix}
\end{PresCodeTexPL}

\begin{PresCodeTexPL}{listing only}
%sudomaths
\SudoMaths[options]{liste}
\pflsudomaths

\begin{EnvSudoMaths}[options]{grille}...\end{EnvSudoMaths}
\end{PresCodeTexPL}

\begin{PresCodeTexPL}{listing only}
%fractales
\FractaleTikz(*)[clés]<options tikz>
\pflfractaltikz

\EtapesFloconKoch[clés]{étapes}
\EtapesTapisSierpinski[clés]{étapes}
\end{PresCodeTexPL}

\begin{PresCodeTexPL}{listing only}
%chateau de cartes
\ChateauCartes[clés]{niveau}<options tikz>
\pflchateaucartes
\end{PresCodeTexPL}

\begin{PresCodeTexPL}{listing only}
%allumettes, dans un environnement tikz
\PfLAllumettes[clés]{liste départs>arrivées}
\pflallumettes
\end{PresCodeTexPL}

\begin{PresCodeTexPL}{listing only}
%commandes sur les dates
\JourSelonDate(*){jj/mm/aaaa}
\pfljourselondate

\DateComplete(*){jj/mm/aaaa}
\pfldatecomplet

\NbJoursEntreDates(*){jj/mm/aaaa}{jj/mm/aaaa}[\macro]
\pflnbjoursdates

\AjoutJoursDate(*){jj/mm/aaaa}{nb}
\pflajoutjoursdate
\end{PresCodeTexPL}

\begin{PresCodeTexPL}{listing only}
%machine à transformer
\MachineTransformer[clés]{e/s}<options tikz>
\pflmachtransf
\end{PresCodeTexPL}

\begin{PresCodeTexPL}{listing only}
%insertion de clipart, en mode inline ou non
\InsererClipart*[options includegraphics]{clipart}
\pflclipart
\end{PresCodeTexPL}

\begin{PresCodeTexPL}{listing only}
%logo sans calculatrice
\pflnocalc(*)[couleur]
\end{PresCodeTexPL}